\chapter{Enrico Bombieri}\label{ch1974enricobombieri:enrico-bombieri-fields-medal-1974-a-comprehensive-report-on-his-work-and-impact}

\begin{center}\rule{0.5\linewidth}{0.5pt}\end{center}

\textbf{Prepared:} July 2026 \textbf{Subject:} Enrico Bombieri (b.\
November 26, 1940, Milan, Italy) \textbf{Position:} Professor,
Institute for Advanced Study, Princeton \textbf{Doctoral Advisor:}
Giovanni Ricci (University of Milan) \textbf{Fields Medal Citation:}
\emph{``For his major contributions to the theory of prime numbers, to
the theory of minimal surfaces, and to the theory of functions of several
complex variables.''}

\begin{center}\rule{0.5\linewidth}{0.5pt}\end{center}

\begin{prereqs} Analytic number theory (prime number theorem, Dirichlet
L-functions, sieve methods), basic algebraic geometry (algebraic
surfaces, Enriques surfaces), differential geometry of minimal surfaces
(Plateau's problem, geometric measure theory). \end{prereqs}

\begin{objectives} Explain the Bombieri--Vinogradov theorem and its
significance as an unconditional substitute for GRH on average; describe
the large sieve inequality and its role in analytic number theory;
understand how Bombieri, De Giorgi, and Giusti settled the Bernstein
conjecture by constructing a non-flat minimal graph in dimension 8.
\end{objectives}

\section{Problem History}

Bombieri's work spans analytic number theory, differential geometry, and
several complex variables. He received the Fields Medal in 1974 in
Vancouver.

The distribution of primes in arithmetic progressions is a classical
problem. Dirichlet proved $\pi(x; q, a) \sim \mathrm{Li}(x) / \phi(q)$
as $x \to \infty$, but the error term $E(x; q, a)$ is only well
controlled for fixed moduli. Under GRH, one gets
$|E(x; q, a)| \ll x^{1/2} \log^2 x$ for all $q < x$, but GRH is
unproven.

Independently, Bombieri and Vinogradov proved in 1965 that for any
fixed $A > 0$, there exists $B = B(A)$ such that
\[
\sum_{q \le Q} \max_{(a,q)=1} \left| \psi(x; q, a) - \frac{x}{\phi(q)} \right| \ll \frac{x}{(\log x)^A}
\]
where $Q = x^{1/2} / (\log x)^B$---an unconditional average GRH up to
the square-root barrier.

A separate problem arises in differential geometry. Bernstein proved in
1914 that the only entire minimal graph in $\mathbb{R}^3$ is a plane.
The Bernstein problem asks whether this holds in higher dimensions. De
Giorgi (1962), Almgren (1966), and Simons (1968) settled dimensions
$n \le 7$. Simons identified a critical cone in $\mathbb{R}^8$ and
showed it was stable but conjectured it was not area-minimizing. In
1969, Bombieri, De Giorgi, and Giusti proved the Simons cone is
area-minimizing, producing a counterexample in dimension 8.

Bombieri also made deep contributions to several complex variables: his
1970 work on algebraic values of meromorphic maps pioneered multiplier
ideal sheaves. He classified Enriques surfaces in positive
characteristic (1974--1976) with Mumford, discovering new types at
characteristic $p = 2$.

\section{How Was It Solved}

The Bombieri--Vinogradov proof uses three ingredients: Dirichlet
characters, bilinear decompositions, and the large sieve. The
congruence condition is isolated using Dirichlet characters:
\[
\psi(x; q, a) - \frac{x}{\phi(q)} \approx \frac{1}{\phi(q)} \sum_{\chi \bmod q} \bar{\chi}(a)\, \psi(x, \chi).
\]
The von Mangoldt function is decomposed via Vaughan's identity into
Type I and Type II bilinear forms. The large sieve inequality states
that for a trigonometric polynomial $S(x) = \sum_n a_n e(nx)$ and
$\delta$-spaced points $x_1, \dots, x_R$ in $\mathbb{R}/\mathbb{Z}$,
\[
\sum_{r=1}^R |S(x_r)|^2 \le (N + \delta^{-1}) \sum_n |a_n|^2.
\]
Bombieri and Davenport obtained the sharp constant by the duality
principle. Choosing rationals $a/q$ with $q \le Q$ yields the character
form of the sieve.

The Bernstein counterexample reduces the minimal surface equation on
the Simons cone to an ODE in radial coordinates. A phase-plane analysis
shows the existence of a non-flat solution asymptotic to the cone. The
proof that the Simons cone is area-minimizing uses calibration: a
co-oriented foliation of $\mathbb{R}^8 \setminus C$ by stable minimal
hypersurfaces provides a calibrating form.

The Enriques surface classification uses the Galois groups of the
canonical cover: $\mathbb{Z}/2$ (classical), $\mu_2$ (singular), or
$\alpha_2$ (supersingular) in characteristic $2$.

\section{Impact of This Discovery}

The Bombieri--Vinogradov theorem underpins the work on bounded gaps
between primes. The GPY method (2005) and Zhang's 2013 breakthrough
proving $p_{n+1} - p_n < 7 \times 10^7$ infinitely often both relied
on this theorem. Maynard and Polymath8 reduced the bound to $246$.

The Bernstein conjecture resolution showed a sharp phase transition at
dimension $n = 7$, opening new directions in geometric measure theory.
Bombieri's classification of Enriques surfaces \cite{bombieri1973} demonstrated that algebraic
geometry over finite fields harbors phenomena absent in characteristic zero.
The theory of multiplier ideal sheaves has become central in algebraic
geometry (Demailly, Nadel, Siu). The large sieve remains a standard tool
for distributional results for primes and $L$-functions.

\section{Modern Views and Improvements}

The Elliott--Halberstam conjecture would push the Bombieri--Vinogradov
barrier to $Q = x^{1-\epsilon}$, reducing bounded gaps to $6$. Zhang
proved a partial form for smooth moduli, and Maynard's multidimensional
sieve produced further improvements.

In geometric measure theory, De Lellis and Spadaro resolved conjectures
on singular sets of area-minimizing currents. Deligne's Weil conjectures
provide the cohomological framework for zeta functions over finite fields,
complemented by Bombieri's exponential sum bounds. Modern multiplier ideal
sheaves led to the Chen--Donaldson--Sun K\"{a}hler--Einstein metric
existence theorem (2015).

\begin{center}\rule{0.5\linewidth}{0.5pt}\end{center}

\section*{Summary}
\addcontentsline{toc}{section}{Summary}

Bombieri proved the Bombieri--Vinogradov theorem (an unconditional average
GRH), settled the Bernstein conjecture in dimensions $\ge 8$, and classified
Enriques surfaces in all characteristics.

\begin{center}\rule{0.5\linewidth}{0.5pt}\end{center}
